% Part: incompleteness
% Chapter: incompleteness-provability
% Section: lob-thm

\documentclass[../../../include/open-logic-section]{subfiles}

\begin{document}

\olfileid{inc}{inp}{tar}

\olsection{د رښتياوالي نه تعريفېدنه}

د \emph{تعريفېدنې} مفهوم د حساب د ژبې پر صوري معناپوهنې ولاړ دے۔
موږ د حساب په ژبه کښې د فورمولونو او جملو يو سټ بيان کړے دے۔
«مراد تفسير» دا دے چې دغه جملې د طبيعي عددونو په باب دعوې وبولو؛
داسې دعوه رښتيا يا ناسمه کېدے شي۔ $\Struct{N}$ دې هغه
!!{structure} وي چې تعريف‌سيمه ئې~$\Nat$ او د حساب د ژبې
نښو تفسير ئې معياري دے۔ بيا~$\Sat{N}{!A}$ معنا لري: «$!A$
په معياري تفسير کښې رښتيا ده»۔

\begin{defn}
د طبيعي عددونو اړيکه~$R(x_1,\dots,x_k)$ په~$\Struct{N}$ کښې
هله او يوازې هله \emph{تعريفېدونکې} ده چې د حساب د ژبې داسې
فورمول~$!A(x_1,\dots,x_k)$ وي چې د ټولو~$n_1,\dots,n_k$ دپاره
$R(n_1,\dots,n_k)$ هله او يوازې هله برقرار وي چې
$\Sat{N}{!A(\num n_1,\dots,\num
  n_k)}$۔
\end{defn}

په بل ډول، يوه اړيکه په~$\Struct{N}$ کښې هله او يوازې هله
تعريفېدونکې ده چې په تيورۍ~$\Th{TA}$ کښې تمثيلېدونکې وي؛ دلته
$\Th{TA} =
\Setabs{!A}{\Sat{N}{!A}}$ د حساب د رښتينو جملو سټ دے۔ (که دا
خبره سملاسي روښانه نۀ وي، تعريفونو ته بېرته ورشئ او ځان پرې
قانع کړئ۔)

\begin{lem}
هره محاسبه کېدونکې اړيکه په~$\Struct{N}$ کښې تعريفېدونکې ده۔
\end{lem}

\begin{proof}
په آسانۍ کتل کېدے شي چې په~$\Th{Q}$ کښې د يوې اړيکې
تمثيلوونکے فورمول په~$\Struct{N}$ کښې هماغه اړيکه تعريفوي۔
\end{proof}

اوس پوښتنه دا ده چې آيا برعکس هم رښتيا دے؟ يعنې آيا په~$\Struct{N}$
کښې هره تعريفېدونکې اړيکه محاسبه کېدونکې ده؟ ځواب نۀ دے۔ مثلاً:

\begin{lem}
د درېدنې اړيکه په~$\Struct{N}$ کښې تعريفېدونکې ده۔
\end{lem}

\begin{proof}
ياد ولرئ چې د کلين د عادي بڼې قضيه وايي: هره جزوي محاسبه
کېدونکې تابعه~$f$ داسې شاخص~$e$ لري چې د ټولو~$x \in \Nat$
دپاره $f(x) =
U(\umin{s}{T(e,x,s)})$ وي؛ دلته~$U$ او~$T$ ابتدايي بازګشتي
او ځکه کلي دي۔ نو~$f(x)$ هله او يوازې هله تعريف شوې وي (يعنې
محاسبه ودرېږي) چې داسې~$s$ وي چې~$T(e,x,s)$ برقرار وي۔

اوس~$H$ د درېدنې اړيکه واخلو، يعنې
\[
H = \Setabs{\tuple{e,x}}{\lexists[s][T(e, x, s)]}.
\]
$!D_T$ دې په~$\Struct{N}$ کښې~$T$ تعريف کړي۔ نو
\[
H = \Setabs{\tuple{e,x}}{\Sat{N}{\lexists[s][!D_T(\num e, \num x, s)]}},
\]
او په دې توګه~$\lexists[s][!D_T(z, x, s)]$ په~$\Struct{N}$
کښې~$H$ تعريفوي۔
\end{proof}

\begin{prob}
وښيئ چې $Q(n) \defiff n \in \Setabs{\Gn{!A}}{\Th{Q} \Proves !A}$
  په حساب کښې تعريفېدونکې ده۔
\end{prob}

خو پخپله~$\Th{TA}$ څنګه ده؟ آيا په حساب کښې تعريفېدونکې ده؟
يعنې آيا سټ~$\Setabs{\Gn{!A}}{\Sat{N}{!A}}$ په حساب کښې
تعريفېدونکے دے؟ د تارسکي قضيه منفي ځواب ورکوي۔

\begin{thm}
\ollabel{thm:tarski}
د حساب د رښتينو !!{sentence}s سټ په حساب کښې نۀ شي تعريفېدے۔
\end{thm}

\begin{proof}
فرض کړئ~$!D(x)$ دا سټ تعريفوي؛ يعنې~$\Sat{N}{!A}$ هله او
يوازې هله چې~$\Sat{N}{!D(\gn{!A})}$۔ د ثابت ټکي د لمې له مخې
داسې فورمول~$!A$ شته چې
$\Th{Q} \Proves !A \liff \lnot !D(\gn{!A})$، او ځکه
$\Sat{N}{!A \liff \lnot !D(\gn{!A})}$۔ خو بيا~$\Sat{N}{!A}$
هله او يوازې هله چې~$\Sat{N}{\lnot !D(\gn{!A})}$؛ دا له دې
فرض سره ټکر لري چې~$!D(y)$ د حساب د رښتينو جملو سټ تعريفوي۔
\end{proof}

تارسکي دا څېړنه د رښتياوالي پر يوه لا عمومي فلسفي مفهوم پلې کړه۔
د هرې ژبې~$L$ دپاره ئې استدلال وکړ چې د~$L$ د رښتياوالي مناسب
مفهوم بايد د هرې جملې~$X$ دپاره دا شرط پوره کړي:
\begin{quote}
«$X$» هله او يوازې هله رښتيا ده چې $X$۔
\end{quote}
د انګرېزۍ ژبې دپاره د تارسکي ډېره نقل شوې بېلګه دا جمله ده:
\begin{quote}
«واوره سپينه ده» هله او يوازې هله رښتيا ده چې واوره سپينه ده۔
\end{quote}
خو د هرې هغې ژبې دپاره چې د قطري تابع د تمثيل توان ولري، او د
هر ژبني پريديکات~$T(x)$ دپاره، داسې جمله~$X$ جوړولے شو چې وايي:
«$X$ هله او يوازې هله چې $T(\text{«$X$»})$ نۀ وي»۔ که موږ نۀ
غواړو چې د رښتياوالي پريديکات يوه جمله په يو وخت کښې هم رښتيا او
هم ناسمه وبولي، نو د تارسکي پايله دا وه چې د يوې ژبې د ټولو جملو
دپاره د رښتياوالي پريديکات د ژبې له پولو څخه له وتلو پرته نۀ شو
ټاکلے۔ په نورو ټکو، د يوې ژبې د رښتياوالي پريديکات په هماغه ژبه
کښې نۀ شي تعريفېدے۔

\end{document}
